% -*- coding: utf-8 -*-
\documentclass[12pt]{ctexart}

\usepackage[
   paperwidth=285mm,   % 宽度
  paperheight=420mm,  % 高度
    top=1in,
    landscape,
    bottom=1in,
    left=1in,
    right=1in,
    headheight=25pt % Increased height to accommodate the header
]{geometry}

\usepackage{amsmath}
\usepackage{amssymb}
\usepackage{fancyhdr}
\usepackage{multicol}
\usepackage{lipsum}
\usepackage{tikz}
\usepackage{enumitem}
\usepackage{graphicx}
\usepackage{hyperref}
\usepackage{amsthm}
\newenvironment{solution}{\begin{proof}[\indent\bf 解]}{\end{proof}}

% Initialize the fancy header style
\pagestyle{fancy}
\fancyhf{} % Clear all default header and footer fields

% Set the header content
\fancyhead[L]{姓名：} % Left header
%\fancyhead[C]{\textbf{Midterm Examination}} % Center header
%\fancyhead[R]{2025 秋} % Right header

% Set the footer content
%\fancyfoot[L]{Prof. Alan Turing} % Left footer
\fancyfoot[C]{\thepage\ of \pageref{LastPage}} % Center footer (Page X of Y)
%\fancyfoot[R]{Total Marks: 100} % Right footer

% Draw a line below the header
\renewcommand{\headrulewidth}{0.4pt}
% Draw a line above the footer
\renewcommand{\footrulewidth}{0.4pt}


\begin{document}

\begin{center}
    % {\Large \textbf{University of Excellence}}
    % \vspace{2mm}
    % {\large Department of Computer Science}
    % \vspace{8mm}
    {\Huge \textbf{群：通向几何之路 chapter 1}}
%    \vspace{4mm}
    % {\Large CS 101: Introduction to Programming}
    % \vspace{12mm}
\end{center}

% \begin{center}
% \begin{tabular}{l@{\extracolsep{8em}}l}
%     \textbf{Date:} September 2, 2025 & \textbf{Duration:} 3 Hours \\
%     \textbf{Total Marks:} 100 & \textbf{Sections:} 3 \\
% \end{tabular}
% \end{center}

% \vspace{5mm}

% \begin{center}
%     \fbox{\begin{minipage}{0.9\textwidth}
%         \textbf{\Large Instructions:}
%         \begin{itemize}[leftmargin=*]
%             \item Answer all questions in the space provided.
%             \item Write your name and student ID on the top of every page (if required).
%             \item No electronic devices are permitted, except for non-programmable calculators.
%             \item The total marks for this paper is 100. Marks for each question are indicated in brackets.
%             \item Show all your work for full credit.
%         \end{itemize}
%     \end{minipage}}
% \end{center}

% \vspace{10mm}
% \hrule
% \vspace{10mm}


\begin{multicols}{2}

\section*{映射 \(N \rightarrow  L\)}

先观察以下图示。（图例：\(N =\) 数字 \(L =\) 字母）

% \begin{figure}[h]
%   \centering
%   \includegraphics[]{images/fig-1.png}
%   \includegraphics[width = 0.8\linewidth]{images/fig-1.png}
%   \caption{\(f = \left\{ (1,a), (2,a), (3,b), (4,c), (5,c) \right\}\)}
%   \label{fig:1}
% \end{figure}


映射 \(f = \left\{ (1,a), (2,a), (3,b), (4,c), (5,c) \right\}\)
\begin{center}
   \includegraphics[width = 0.8\linewidth]{images/fig-1.png}
\end{center}

映射 \(g = \{ \left( {1,a}\right) ,\left( {2,a}\right) ,\left( {3,a}\right) ,\left( {4,b}\right) ,\left( {5,b}\right) \}\)
% \begin{figure}[h]
%   \centering
%   \includegraphics[width = 0.8\linewidth]{images/fig-2.png}
%   \caption{ }
%   \label{fig:2}
% \end{figure}

\begin{center}
   \includegraphics[width = 0.8\linewidth]{images/fig-2.png}
\end{center}

其中定义域 \(N = \{ 1,2,3,4,5\}\) ，陪域 \(L = \{ a,b,c\}\) 。

\(f\) 和 \(g\) 两者都是从 \(N\) 到 \(L\) 的映射示例：
\[
f : N \rightarrow  L, \quad g : N \rightarrow  L
\]

再观察以下图示。以下皆非  \(N \rightarrow  L\) 的映射。

\begin{center}
  \includegraphics[width = 0.4\linewidth]{images/fig-3.png}
  \includegraphics[width = 0.4\linewidth]{images/fig-4.png}  
\end{center}

\begin{enumerate}
\item 依上述图示及其隐含规则完成下列句子。

  当集合 \(N\) 中的每个元素 \(n\) 都满足 \underline{\hbox to 40mm{}} 时，映射 \(f : N \rightarrow  L\) 即被定义。

\item 依上述图示及其隐含规则完成下列句子。

  映射图像所用的矩形阵列表示所谓的笛卡尔积 \(N \times  L\) 。 \(N \times  L\) 的每个元素都是有序对$(n, x)$，其中 \(n\) 来自 \(N\) ， \(x\) 来自 \(L\) 。映射 \(N \rightarrow  L\) 的图像由 \(N \times  L\) 的一个子集构成，对 \underline{\hbox to 40mm{}} 恰好只有一个 $(n, x)$ 与之对应。

\item 下列哪些定义了实数 \(\mathbb{Q} \rightarrow  \mathbb{Q}\) 的映射？
\begin{enumerate}
\item \(x \mapsto  {x}^{2}\) ,
\item  \({x}^{2} \mapsto  x\) ，
\item  \(x \mapsto  1/x\) ，
\end{enumerate}

\item  若 \(A = \{ 0,1\}\) ，存在多少种不同的映射 \(A \rightarrow  A\) ？

\end{enumerate}





\section*{单射}

\begin{enumerate}
  \setcounter{enumi}{4}
\item 绘制下列映射 \(\mathbb{Q} \rightarrow  \mathbb{Q}\) 的示意图
\begin{enumerate}
\item  \(x \mapsto  {x}^{3}\) 
\item  \(x \mapsto  {x}^{2}\) 。
\end{enumerate}
\item 若 \({x}^{3} = {y}^{3}\) ，是否可推出 \(x = y\) ？若 \({x}^{2} = {y}^{2}\) ，是否可推出 \(x = y\) ？

这里的区别使我们称由 \(\alpha  : x \mapsto  {x}^{3}\) 定义的 \(\alpha  : \mathbb{R} \rightarrow  \mathbb{R}\) 为单射。我们说 \(\beta  : x \mapsto  {x}^{2}\) 在 \(\mathbb{R}\) 上不是单射的。

\item 设 \(A = \{ 1,2,3,4\}\) 。
  \begin{enumerate}
  \item 展示一个单射 \(A \rightarrow  A\) 的图像。
  \item 展示一个非单射 \(A \rightarrow  A\) 的图像。
  \item 存在多少个映射 \(A \rightarrow  A\) ？
  \end{enumerate}
\item 设 \(\mathbb{N}\) 表示自然数集。绘制由 \(n \mapsto  {n}^{2}\) 定义的映射 \(\mathbb{N} \rightarrow  \mathbb{N}\) 的部分图像。这是单射吗？
\item 若已知某映射是单射，其图像的行有何特征？
\end{enumerate}


\section*{满射}

\begin{enumerate}
  \setcounter{enumi}{8}
\item 绘制下列映射 \(\mathbb{Q} \rightarrow  \mathbb{Q}\) 的示意图:
\begin{enumerate}
\item \(\alpha  : x \mapsto  {2}^{x}\) 
\item \(\beta  : x \mapsto  x + 1\) 。
\end{enumerate}
对于任意 \(y \in \mathbb{Q}\) ，是否总能找到 \(x \in \mathbb{Q}\) 使得 \(\alpha  : x \mapsto  y\) 成立？

对于任意实数 \(y \in \mathbb{Q}\) ，是否总能找到 \(x\in \mathbb{Q}\) 使得 \(\beta  : x \mapsto  y\) 成立？

由此差异可称 \(\beta\) 为满射，而 \(\alpha\) 不是满射。

\item  设 \(A = \{ 1,2,3,4\}\) 。展示一个映射 \(A \rightarrow  A\) 的图像，使得其分别满足
\begin{enumerate}
\item 该映射是满射。
\item 该映射不是满射。
\end{enumerate}


\item  绘制由以下定义映射 \(\mathbb{N} \rightarrow  \mathbb{N}\) 的部分图像:
  $$
n \mapsto
\begin{cases}
  \frac{n}{2} & n \text{为偶数} \\
  \frac{n+1}{2} & n \text{为奇数}
\end{cases}
  $$
这是单射吗？这是一个满射吗？

\item  关于满射映射图像的行，我们能得出什么结论？

\item  设 \(A = \{ 1,2,3,4\}\) 。
\begin{enumerate}
\item 能否构造一个单射但不满射的映射 \(A \rightarrow  A\) ？
\item 能否构造一个满射但不单射的映射 \(A \rightarrow  A\) ？
\end{enumerate}

\item 能否构造一个单射但不满射的映射 \(\mathbb{N} \rightarrow  \mathbb{N}\) ？能否构造一个满射但不单射的映射 \(\mathbb{N} \rightarrow  \mathbb{N}\) ？

\item  推测任意集合 \(A\) 所需满足的条件，使得每个单射映射 \(A \rightarrow  A\) 都必须是满射的，且每个满射映射 \(A \rightarrow  A\) 都必须是单射的。借助问题8和问题12证明你的猜想。

既是单射又是满射的映射称为双射。在定义域和陪域相同的有限集映射中，单射、满射和双射实际上无法区分。

\item  举例说明映射 \(\mathbb{Q} \rightarrow  \mathbb{Q}\) 的以下情形:
\begin{enumerate}
\item 单射且满射，

\item  单射但非满射，

\item  满射但非单射，

\item  既非单射也非满射。
\end{enumerate}
 
\item 若存在一个单射但不满射的映射 \(A \rightarrow  A\) ，关于集合 \(A\) 能得出什么结论？

\end{enumerate}

\section*{映射的复合}

\begin{enumerate}
  \setcounter{enumi}{17}
\item 若 \(\alpha\) 和 \(\beta\) 是由以下定义确定的映射 \(\mathbb{Q} \rightarrow  \mathbb{Q}\)

\[\alpha  : x \mapsto  {2x}, \quad \beta  : x \mapsto  x + 1\]

那么我们定义 \({x\beta } : x \mapsto  {2x} + 1\) 。

\[
x \mapsto  {2x} \mapsto  {2x} + 1
\]

我们将其记为 \(\left( x\right) {\alpha \beta } = \left( {2x}\right) \beta  = {2x} + 1\) 。

根据类似定义确定 \(\left( x\right) {\beta \alpha}\) 。

\item  若 \(\alpha  : A \rightarrow  B\) 和 \(\beta  : B \rightarrow  C\) 是映射，请通过确定 \(\left( x\right) {\alpha \beta }\) 给出 \({\alpha \beta } : A \rightarrow  C\) 的形式化定义。

映射 \({\alpha \beta }\) 称为映射 \(\alpha\) 与 \(\beta\) 的复合。

\item 已知 $\alpha, \beta$ 的示意图如下，试绘制 $\alpha\beta$ 的示意图。

  \begin{center}
     \includegraphics[width = 0.4\linewidth]{images/fig-5.png}
  \end{center}

\item   “\(\alpha  : A \rightarrow  B\) 是单射”，试给出其形式定义。

\item  若 \(\alpha  : A \rightarrow  B\) 和 \(\beta  : B \rightarrow  C\) 都是单射，能推出 \({\alpha \beta } : A \rightarrow  C\) 的什么性质？

\item  设 \(A = \{ 1,2,3\}\) 和 \(B = \{ 1,2,3,4\}\) ，并令 \(\alpha  : A \rightarrow  B\) 为单射。

能否构造满足以下条件的映射 \(\beta  : B \rightarrow  A\) :

\begin{enumerate}
\item \({\beta \alpha }\) 是恒等映射 \(B \rightarrow  B\) ，
\item  \({\alpha \beta }\) 是恒等映射 \(A \rightarrow  A\) 
\end{enumerate}
（恒等映射是将每个元素都映到自身的映射）
\end{enumerate}


\section*{逆映射}
\begin{enumerate}
  \setcounter{enumi}{23}
\item 设 \(A\) 和 \(B\) 为任意集合，且 \(\alpha  : A \rightarrow  B\) 为单射。请说明如何定义 \(\beta  : B \rightarrow  A\) 使得 \({\alpha \beta }\) 成为 \(A\) 上的恒等映射。此时映射 \(\beta\) 称为 \(\alpha\) 的右逆。因此单射具有右逆。
\item  “\(\alpha  : A \rightarrow  B\) 是满射”，试给出其形式定义。
\item 若 \(\alpha  : A \rightarrow  B\) 和 \(\beta  : B \rightarrow  C\) 都是满射，复合映射 \({\alpha \beta } : A \rightarrow  C\) 具有什么性质？
\item 设 \(A = \{ 1,2,3,4\}\) 和 \(B = \{ 1,2,3\}\) ，且 \(\alpha  : A \rightarrow  B\) 为满射。
\begin{enumerate}
\item 能否构造映射 \(\beta  : B \rightarrow  A\) 使得 \({\alpha \beta }\) 成为 \(A\) 上的恒等映射？
\item 能否构造映射 \(\beta  : B \rightarrow  A\) 使得 \({\beta \alpha } : B \rightarrow  B\) 成为 \(B\) 上的恒等映射？
\end{enumerate}

\item 设 \(A\) 和 \(B\) 为任意集合，其中 \(\alpha  : A \rightarrow  B\) 是满射。证明如何定义映射 \(\beta  : B \rightarrow  A\) 使得 \({\beta \alpha }\) 成为 \(B\) 上的恒等映射？

  此时映射 \(\beta\) 称为 \(\alpha\) 的左逆。故满射必存在左逆。
\item 设 \(A = \{ 1,2,3\}\) 。列出所有双射 \(A \rightarrow  A\) ，即 \(A\) 的所有排列。求每个双射的左逆和右逆。
\item 对任意双射 \(\alpha  : A \rightarrow  B\) ，定义双射 \(\beta  : B \rightarrow  A\) 使得 \({\alpha \beta }\) 为恒等映射 \(I_A : A \rightarrow  A\) 且 \({\beta \alpha }\) 为恒等映射 \(I_B : B \rightarrow  B\) 。

证明 \({\alpha \beta } = I_A : A \rightarrow  A\) 或 \({\beta \alpha } = I_B : B \rightarrow  B\) 能唯一确定 \(\beta\) 。

因此双射具有双侧逆。
\end{enumerate}


\section*{封闭性}

\begin{enumerate}
\setcounter{enumi}{30}
\item 若 \(\alpha\) 和 \(\beta\) 都是双射 \(A \rightarrow  A\) ，则 \({\alpha \beta } : A \rightarrow  A\) 具有什么性质？
\end{enumerate}

\section*{结合律}

\begin{enumerate}
  \setcounter{enumi}{31}
\item 设 \(\alpha  : A \rightarrow  B,\beta  : B \rightarrow  C\) 和 \(\gamma  : C \rightarrow  D\) 为映射。利用问题19的定义证明:对于 \(A\) 的每个元素 \(x\) ，都有

\[\left( x\right) \left\lbrack  {\left( {\alpha \beta }\right) \gamma }\right\rbrack   = \left( x\right) \left\lbrack  {\alpha \left( {\beta \gamma }\right) }\right\rbrack  ,\]

由此，两个映射 \(A \rightarrow  D,\left( {\alpha \beta }\right) \gamma\) 与 \(\alpha \left( {\beta \gamma }\right)\) 不可区分。

定理 \(\left( {\alpha \beta }\right) \gamma  = \alpha \left( {\beta \gamma }\right)\) 称为映射复合运算的结合律。

\item  运用结合律证明:对于四个映射 \(\alpha  : A \rightarrow  B\) 、 \(\beta  : B \rightarrow  C,\gamma  : C \rightarrow  D,\delta  : D \rightarrow  E\) ，有 \(\alpha \left\lbrack  {\beta \left( {\gamma \delta }\right) }\right\rbrack   = \left\lbrack  {\left( {\alpha \beta }\right) \gamma }\right\rbrack  \delta\) 。

  \section*{对称群}

\item  设 \(\mathfrak{S}\) 为所有双射 \(A \rightarrow  A\) 的集合。证明下列每个命题。
\begin{enumerate}
\item 【封闭性】若 \(\alpha\) 和 \(\beta\) 属于 \(\mathfrak{S}\) ，则 \({\alpha \beta }\) 也属于 \(\mathfrak{S}\) 。

\item 【结合性】若 \(\alpha ,\beta\) 和 \(\gamma\) 属于 \(\mathfrak{S}\) ，则 \(\alpha \left( {\beta \gamma }\right)  = \left( {\alpha \beta }\right) \gamma\) 。

\item  【单位元】由 \(I : x \mapsto  x\) 定义的 \(I : A \rightarrow  A\) (对所有 \(A\) 中的 \(x\) )属于 \(\mathfrak{S}\) 。对所有 \(\mathfrak{S}\) 中的 \(\alpha\) 有 \({I\alpha } = {\alpha I} = \alpha\) 。

\item  【逆元】对于 \(\mathfrak{S}\) 中的每个 \(\alpha\) ，存在 \(\mathfrak{S}\) 中的 \(\beta\) 使得 \({\alpha \beta } = {\beta \alpha } = I\) 。
\end{enumerate}
这四个定理可总结为:集合到自身的双射在复合运算下构成群。此情形下的群称为 \(A\) 上的对称群，记作 \({\mathfrak{S}}_{A}\) 。
\end{enumerate}






\end{multicols} % End the two-column environment


% Use \newpage or \clearpage to start a new page if needed.
% \clearpage


% Label this page as 'LastPage' for the footer reference
\label{LastPage}
\end{document}
% =============================================================================
% END OF DOCUMENT
% =============================================================================

%%% Local Variables:
%%% mode: latex
%%% TeX-master: t
%%% End:
